%
% dial_settlement.tex
%
% A Z specification of how a dial is ENDED and SETTLED: the reader slot, the
% listener's start and stop, the inline reader's last check and release, the
% no-reader settle, the Hub's own end and its bound.  It refines the two
% abstractions that docs/crash_attribution.tex makes (a "reader slot" that
% someone holds, a "settle" that someone performs) to the granularity of the
% code, because the final review rounds on the lux-ghhf branch found, again and
% again, a path on which an ended dial is never settled or is settled undrained
% (R7-F1, R8 R1, and the three defects of round 9).  This is the recurrence rule:
% the same class of defect in two or more rounds is a state-space problem, and
% the state space is finite.
%
% THIS FILE IS A FIDELITY CONTROL: dial_settlement.tex with the release and the
% debt read in two steps.  S1a violated.  The controls, in docs/:
%   dial_settlement_start_fail_buggy.tex       start_on on an ended dial keeps the slot
%   dial_settlement_stop_skip_drain_buggy.tex  the listener's exit reads ended twice
%                                              (a substitution of its steps, not a
%                                              one-point mutation)
%   dial_settlement_inline_strand_buggy.tex    a refused settle leaves no debt
% Also in docs/:
%   dial_settlement_split_settle_buggy.tex     the post-drain record is not first-wins
%   dial_settlement_atomic_refuse_buggy.tex    refusal and debt are two steps
%   dial_settlement_atomic_release_buggy.tex   release and debt-read are two steps
%   dial_settlement_defer_only_buggy.tex       a Hub request with no bound (T15 absent)
%   dial_settlement_timer_hub_only_buggy.tex   only a Hub request arms the bound
%   dial_settlement_s3_unguarded_buggy.tex     a capture without the current guard
%   dial_settlement_s4_overwrite_buggy.tex     a recorder that overwrites
%   dial_settlement_hub_record_buggy.tex       the Hub end records hub at once (S5)
%   dial_settlement_hub_peek_buggy.tex         the peek rule, an ack ahead of the FIN
% The code it governs: src/punt_lux/domain/hub/{display_connection,dial_reader,
% display_listener,listener_thread,display_replies,reply_sink,dial_generation,
% reader_slot,dial_drain}.py.  Companions: crash_attribution.tex (what a capture
% charges), listener_locks.tex (the lock order rq < sl < st < lk < qu).
%
% Type-check:  fuzz docs/dial_settlement.tex
% Model-check: see the section Verification (one run at a time; each is a few
%              seconds and well under a gigabyte).
%
\documentclass[a4paper,10pt,fleqn]{article}
\usepackage[margin=1in]{geometry}
\usepackage{fuzz}

\begin{document}

\title{Dial Settlement: a Z Specification}
\author{Who ends a dial, who settles it, and who may let go}
\date{October 2026}
\maketitle

\section{What is modelled}

A \emph{dial} is one generation of the Hub's connection to the display.  It is
\emph{published} (current) from its birth until some one party \emph{captures}
it: takes its ledger, queues its death or carries its suspects, and makes it
not current.  Between the first end and the capture the dial is ended and
unsettled.  Five properties are claimed of that interval.

\begin{description}
  \item[S1] No dial stays ended (or end-requested), current and unsettled for
    ever.  Two forms.  \emph{S1a}, with a holder that makes progress: an end is
    always followed by a settle, by a party that was already running, with no
    timer.  \emph{S1b}, with a holder that never makes progress: the bound of
    rule~6 settles it.  A dial that is wedged is one that redial cannot replace,
    because it is still current.
  \item[S2] A dial ended by a \emph{dead} display is captured only after a drain
    that began after the end was recorded.  (The model's \texttt{cdisp} is
    always a dead display; see below for the live display and the carve-outs.)
  \item[S3] Exactly one capture per dial.  \emph{Structural:} every capture is
    guarded by \texttt{current} and clears it, so the flag cannot be set; the
    content of S3 is the guard, and the control \texttt{s3\_unguarded} removes it
    and trips the flag.
  \item[S4] The cause that is captured is the first one recorded.
    \emph{Structural:} \texttt{ended} has one first-wins writer; the control
    \texttt{s4\_overwrite} lets a recorder overwrite and trips the flag.
  \item[S5] A display that was already dead when the settler observed the socket
    is never recorded as a Hub end.  \emph{Structural in the design} (the cause
    is a function of the observation); the controls \texttt{hub\_record} and
    \texttt{hub\_peek} trip it.
\end{description}

S1 is checked as reachability of a wedge state and by the deadlock check; S2 to
S5 are flags set by the operation that would violate them, and the goal is a
state with a flag set.  The reason is the one \texttt{listener\_locks.tex}
gives: ProB treats the primed state schema as part of each operation's guard, so
a violation written as a conjunct disables the operation instead of reporting
it.

\paragraph{The carve-outs of S2, with their bounds.}
Three paths capture a display or unresponsive end without the drain, because the
drain is not possible or not safe.  Each sets its own flag, never \texttt{s2bad}.
\begin{enumerate}
  \item \emph{Decode errors} (\texttt{s2dec}).  A frame that broke the stream
    (the reader's malformed-frame exit, \texttt{end\_stream(display)} in an
    inline read) cannot be read past.  The acknowledgements behind the bad frame
    are unread, and their scenes may be charged.  Bound: the acks in flight on the
    one dial, once.
  \item \emph{The bound of rule~6} (\texttt{s2blk}).  The Hub, tired of waiting
    for a holder that has not settled, captures without draining, whether the
    recorded cause is \texttt{display} or \texttt{unresponsive}.  Bound: every
    acknowledgement queued while the holder stalls, including what sits
    undelivered in the holder's frame buffer, and only after the bound (the stop
    join, two seconds).  A bound capture of a display end may therefore charge
    scenes whose acknowledgements were never read.  That is an accepted
    over-charge (DES-102).
  \item \emph{A live display} (unresponsive ends, for the holder's settle).  A
    display that is alive may send an acknowledgement after any drain, so no
    drain is ever complete.  S2 for an unresponsive end, settled by a holder, is
    that a drain that began after the record precedes the capture; an ack that
    arrives after that drain is bounded by one drain-to-capture gap.  This holds
    for the holder's settle only; the bound's capture of an unresponsive end is
    carve-out~2.  A dead display has its end of file queued behind its last byte
    on a same-host socket, so its drain is complete there; the claim that a drain
    is complete is restricted to same-host sockets, and that restriction is the
    contract (see \emph{What is not modelled}).  The classification of a
    send timeout as the display's end (a live display charged as dead) is a
    separate bound, B3, and is not modelled here.
\end{enumerate}

\section{The design}

Settlement belongs to the \emph{holder of the reader slot}; everyone else only
\emph{asks}.

\begin{enumerate}
  \item \textbf{Capture authority is the slot.}  Only a holder of the slot
    (the listener thread, an inline reader, or a releaser taking over) captures a
    dial.  In code: \texttt{close\_if\_current} leaves the connection's
    interface; the slot hands its holder a \texttt{Settler}, the one object with
    a \texttt{settle()} method, and the listener, the inline reader and the
    no-reader settle are the only holders.  There is no flag with a default
    (\texttt{by\_reader=True}) and no Hub door.  The model's one operation that
    captures without holding the slot is the bound of rule~6, and it is
    labelled.
  \item \textbf{Ends record or request; they never capture.}  A display or
    unresponsive end records its cause, first wins.  A Hub end sets a request
    flag (\texttt{hubreq}) and records no cause.  Both then \emph{ask for a
    settle}: take the slot if it is free and settle at once; if it is held, leave
    a debt on the slot and wait for the bound.
  \item \textbf{Whoever lets go of the slot honours the debt.}  A release
    (listener start failure, inline exit) that finds a debt keeps the slot,
    settles, and looks again.  Every path that gives up the reader role releases
    the slot this way, including a start that raises because the dial was ended.
  \item \textbf{Settle is two steps, a drain and a decision.}  The drain reads
    the socket to the point where it would block and notes whether end of file
    came.  The decision fixes the cause with a first-wins record: the observed
    end of file is \texttt{display}, else a pending Hub request is \texttt{hub};
    a cause recorded earlier stands.  If the record stands and was made after the
    drain began, the drain was not complete: \emph{the settler drains again}
    before it captures.  After the record is fixed nothing can change the cause,
    so the second drain is the last.  Both settlers, the listener and the
    releaser, run these steps.
  \item \textbf{A Hub request acts at once on what it can.}  It refuses writes
    and fails waiters with the Hub's error immediately, without fixing a cause.
    It does not capture, and it does not redial.
  \item \textbf{The wait is bounded (T15).}  An asker refused behind a held slot
    waits.  If the holder has not settled within the bound (the stop join,
    2~s), the asker records its cause (a Hub request records \texttt{hub}) with
    the same first-wins record, and captures.  A capture of cause \texttt{hub}
    charges nothing; its scenes are carried.  The holder's later settle finds the
    dial not current and does nothing.  The bound applies to \emph{every} asker,
    not only to a Hub request: a display end discovered behind a holder that
    never looks is otherwise wedged for as long as the holder stalls
    (the control \texttt{timer\_hub\_only}).
\end{enumerate}

Rules 4 and 5 are what close S5.  The model assumes the observation is one
linearization point (the drain to \texttt{EAGAIN} together with the end-of-file
test).  A display that dies after it dies after the Hub's end; that is not a
defect, and no rule can distinguish it.  The residual widens to the bound if
rule~6 fires: the Hub records \texttt{hub} for a display it has not observed, and
the flag \texttt{s5blk} says so.

\paragraph{The three atomicity obligations.}
S1a holds only if each of these is one critical section of the slot's condition
lock; each has a control that splits it and fails.  With the timer present a
broken obligation does not wedge: it degrades to the bound of rule~6, that is to
the two-second wait and an undrained capture (\texttt{s2blk}).
\begin{enumerate}
  \item \textbf{Refusal and debt}: observing the slot held and setting the debt
    (\texttt{atomic\_refuse}).
  \item \textbf{Release and debt-read}: reading the debt and freeing the slot
    (\texttt{atomic\_release}).
  \item \textbf{The post-drain record}: the first-wins record of the decided
    cause, with the re-drain when it is lost (\texttt{split\_settle}).  This one
    is the dial's \texttt{end()} under the sink's own lock, not the slot's.
\end{enumerate}

\section{Carriers and bounds}

One dial.  Every operation in the code is scoped to a generation: a reader, a
slot, a write failure, and a capture all name the dial they act on, and the Hub
redials only when the dial is no longer current.  A second dial therefore runs
in a disjoint product with the first, and cannot interact with it; the model
includes \texttt{Redial}, enabled only on an unpublished dial, to say so.  Five
actors, each of which acts at most once: the starter (\texttt{start\_on}), the
stopper (a Hub end: \texttt{stop} or \texttt{close}), a discoverer of a display
or unresponsive end (a write failure or the keepalive), one inline reader, and
the listener thread.  Socket content is zero or one unread acknowledgement,
which is enough to tell a drained capture from an undrained one; it arrives at
any time while the display lives.  The dial is born holding the listener's slot
(the registry's \texttt{open\_listening}) or free (a library link).  A holder may
\emph{block} (a handler on the listener thread, an inline read that never
returns): it then takes none of its own steps.

Send timeouts that do not abort a wedged \texttt{sendall} are a separate
obligation (B3 and the keepalive), outside this model.

\begin{zed}
  CAUSE ::= nocause | cdisp | cunresp | chub
\end{zed}

\begin{zed}
  SLOT ::= sfree | slistener | sinline
\end{zed}

\begin{zed}
  FLAG ::= fclear | fset
\end{zed}

\begin{zed}
  ACTOR ::= aSt | aSp | aDi | aIn | aRd
\end{zed}

The program counters.  \texttt{pst0}: the starter has not yet started the
listener.  \texttt{prel0}: it failed and owes a release.  \texttt{phub}: a Hub
end is to be requested.  \texttt{pend}: a display end is to be recorded.
\texttt{psu}: the no-reader settle is to be tried.  \texttt{pwait}: refused, and
waiting for the bound.  \texttt{pset}: this actor holds the slot and owes a
drain.  \texttt{pdec}: it has drained and owes the decision.  \texttt{prel}: it
is about to release.  \texttt{iread}, \texttt{ilast}: the inline reader is
reading, and has finished reading and will make its last check.  \texttt{rtop}:
the listener is at the top of its loop.  \texttt{rnone}: no listener exists.

\begin{zed}
  PC ::= pidle | pst0 | prel0 | phub | pend | psu | pwait | pset | pdec | prel |
    prelf | pdone | iread | ilast | rtop | rnone | rdone
\end{zed}

\begin{schema}{Dial}
  ended : CAUSE \\
  slot : SLOT \\
  debt : FLAG \\
  current : FLAG \\
  sock : 0 \upto 1 \\
  capN : 0 \upto 2 \\
  hubreq : FLAG
\end{schema}

\begin{schema}{Env}
  launched : FLAG \\
  eof : FLAG \\
  redialed : FLAG \\
  blocked : FLAG
\end{schema}

\begin{schema}{Flags}
  s2bad : FLAG \\
  s3bad : FLAG \\
  s4bad : FLAG \\
  s5bad : FLAG
\end{schema}

\begin{schema}{Carve}
  s2blk : FLAG \\
  s2dec : FLAG \\
  s5blk : FLAG
\end{schema}

\begin{schema}{Pc}
  pc : ACTOR \fun PC
\end{schema}

\begin{schema}{Dc}
  dcs : CAUSE
\end{schema}

\begin{schema}{RLoc}
  rcause : CAUSE \\
  rdec : FLAG \\
  rdb : FLAG
\end{schema}

\begin{schema}{Loc}
  obs : ACTOR \fun CAUSE \\
  dsaw : ACTOR \fun CAUSE
\end{schema}

\texttt{ended} is the first recorded cause.  \texttt{hubreq} is a request, not a
cause.  \texttt{eof} says that the display has died and its end of file is
queued behind whatever is unread; \texttt{sock} is that unread
acknowledgement.  \texttt{sock} stands for the socket \emph{and} the reader's
undelivered buffer: that is sound only because rule~1 makes the settler the sole
owner of the frame buffer, so what the reader has taken and not delivered is the
settler's own, and a drain by the settler leaves neither (the gluing is in the
section Refinement).  \texttt{dcs} is the discoverer's cause.  \texttt{obs} is
what an actor's last drain saw (end of file, or nothing); \texttt{dsaw} is the
cause that stood when that drain began.  \texttt{rcause}, \texttt{rdec} and
\texttt{rdb} are locals that only controls use.

\begin{schema}{Sys}
  Dial \\ Env \\ Flags \\ Carve \\ Pc \\ Dc \\ RLoc \\ Loc
\end{schema}

\begin{axdef}
  nxt : ACTOR \fun PC
\where
  nxt = \{ aSt \mapsto prel, aSp \mapsto prel, aDi \mapsto prel, aIn \mapsto prel,
    aRd \mapsto rdone \}
\end{axdef}

Where an actor goes after its settle: a caller to its release, the listener to
its end.

\begin{schema}{Init}
  Sys'
\where
  ended' = nocause \\
  debt' = fclear \\
  current' = fset \\
  sock' \in 0 \upto 1 \\
  capN' = 0 \\
  hubreq' = fclear \\
  launched' = fclear \\
  eof' = fclear \\
  redialed' = fclear \\
  blocked' = fclear \\
  s2bad' = fclear \\
  s3bad' = fclear \\
  s4bad' = fclear \\
  s5bad' = fclear \\
  s2blk' = fclear \\
  s2dec' = fclear \\
  s5blk' = fclear \\
  dcs' = cdisp \\
  rcause' = nocause \\
  rdec' = fclear \\
  rdb' = fclear \\
  obs' = ACTOR \cross \{ nocause \} \\
  dsaw' = ACTOR \cross \{ nocause \} \\
  (slot' = slistener \land pc' = \{ aSt \mapsto pst0, aSp \mapsto pidle,
    aDi \mapsto pidle, aIn \mapsto pidle, aRd \mapsto rnone \}) \lor
  (slot' = sfree \land pc' = \{ aSt \mapsto pdone, aSp \mapsto pidle,
    aDi \mapsto pidle, aIn \mapsto pidle, aRd \mapsto rnone \})
\end{schema}

The frames.  Each operation names, through one of these, which parts of the
state it may change; the rest are unchanged.

\begin{schema}{OnlyPc}
  \Xi Dial \\ \Xi Env \\ \Xi Flags \\ \Xi Carve \\ \Delta Pc \\ \Xi Dc \\
  \Xi RLoc \\ \Xi Loc
\end{schema}

\begin{schema}{DialPc}
  \Delta Dial \\ \Xi Env \\ \Xi Flags \\ \Xi Carve \\ \Delta Pc \\ \Xi Dc \\
  \Xi RLoc \\ \Xi Loc
\end{schema}

\begin{schema}{DialOnly}
  \Delta Dial \\ \Xi Env \\ \Xi Flags \\ \Xi Carve \\ \Xi Pc \\ \Xi Dc \\
  \Xi RLoc \\ \Xi Loc
\end{schema}

\begin{schema}{EnvOnly}
  \Xi Dial \\ \Delta Env \\ \Xi Flags \\ \Xi Carve \\ \Xi Pc \\ \Xi Dc \\
  \Xi RLoc \\ \Xi Loc
\end{schema}

\begin{schema}{EnvPc}
  \Xi Dial \\ \Delta Env \\ \Xi Flags \\ \Xi Carve \\ \Delta Pc \\ \Xi Dc \\
  \Xi RLoc \\ \Xi Loc
\end{schema}

\begin{schema}{PcDc}
  \Xi Dial \\ \Xi Env \\ \Xi Flags \\ \Xi Carve \\ \Delta Pc \\ \Delta Dc \\
  \Xi RLoc \\ \Xi Loc
\end{schema}

\begin{schema}{Drain}
  \Delta Dial \\ \Xi Env \\ \Xi Flags \\ \Xi Carve \\ \Delta Pc \\ \Xi Dc \\
  \Xi RLoc \\ \Delta Loc
\end{schema}

\begin{schema}{Cap}
  \Delta Dial \\ \Xi Env \\ \Delta Flags \\ \Xi Carve \\ \Delta Pc \\ \Xi Dc \\
  \Xi RLoc \\ \Xi Loc
\end{schema}

\begin{schema}{CapCarve}
  \Delta Dial \\ \Xi Env \\ \Delta Flags \\ \Delta Carve \\ \Delta Pc \\ \Xi Dc \\
  \Xi RLoc \\ \Xi Loc
\end{schema}

\section{The environment}

\begin{schema}{Die}
  EnvOnly
\where
  eof = fclear \\
  eof' = fset \\
  launched' = launched \\
  redialed' = redialed \\
  blocked' = blocked
\end{schema}

The display dies.  Its end of file is queued behind any acknowledgement it has
already sent.

\begin{schema}{Ack}
  DialOnly
\where
  current = fset \\
  eof = fclear \\
  sock = 0 \\
  sock' = 1 \\
  ended' = ended \\
  slot' = slot \\
  debt' = debt \\
  current' = current \\
  capN' = capN \\
  hubreq' = hubreq
\end{schema}

A living display sends an acknowledgement, at any time.  This is what the round-1
model lacked: without it no capture could ever race an ack.

\begin{schema}{Block}
  EnvOnly
\where
  blocked = fclear \\
  blocked' = fset \\
  launched' = launched \\
  eof' = eof \\
  redialed' = redialed
\end{schema}

The holder of the slot stops making progress, for good: a handler that never
returns.  It takes none of the steps that need it (\texttt{RNotice},
\texttt{IReadIdle}, \texttt{IReadNotice}, \texttt{IBad}, \texttt{RBad}).  Steps
that a holder takes once it has begun to settle are not blocked: the drain does
not wait on the display.

\begin{schema}{Redial}
  EnvOnly
\where
  current = fclear \\
  redialed = fclear \\
  redialed' = fset \\
  launched' = launched \\
  eof' = eof \\
  blocked' = blocked
\end{schema}

The Hub dials again, and only on a dial that is no longer current.  A dial
that is ended but still current is never replaced: that is the wedge.

\begin{schema}{Halt}
  \Xi Sys
\where
  current = fclear \lor (ended = nocause \land hubreq = fclear)
\end{schema}

The model may stop only when the dial is settled or has not been ended.  A
state with the dial ended or end-requested, still current, and nothing enabled
is a deadlock, and probcli reports it.

\section{The starter}

\begin{schema}{StOk}
  EnvPc
\where
  pc~aSt = pst0 \\
  current = fset \\
  ended = nocause \\
  hubreq = fclear \\
  launched' = fset \\
  eof' = eof \\
  redialed' = redialed \\
  blocked' = blocked \\
  pc' = pc \oplus \{ aSt \mapsto pdone, aRd \mapsto rtop \}
\end{schema}

\texttt{start\_on}: the dial is not closed, and the thread starts.

\begin{schema}{StFail}
  OnlyPc
\where
  pc~aSt = pst0 \\
  pc' = pc \oplus \{ aSt \mapsto prel0 \}
\end{schema}

\texttt{start\_on} fails: the dial was found ended or requested
(\texttt{ConnectionClosedError}), or the thread could not be started.  The
design keeps this unguarded, which is the stronger claim: it releases on both.
The starter owes a release of the slot it was born holding.

\begin{schema}{StRel0}
  DialPc
\where
  pc~aSt = prel0 \\
  slot = slistener \\
  ended' = ended \\
  current' = current \\
  sock' = sock \\
  capN' = capN \\
  hubreq' = hubreq \\
  (debt = fset \land debt' = fclear \land slot' = sinline \land
    pc' = pc \oplus \{ aSt \mapsto pset \}) \lor
  (debt = fclear \land debt' = debt \land slot' = sfree \land
    pc' = pc \oplus \{ aSt \mapsto phub \})
\end{schema}

The release.  With a debt the starter keeps the slot and settles; without, it
frees the slot and goes on to request the Hub's end of the dial it could not
start (\texttt{retire(born, hub)}).

\section{The ends}

\begin{schema}{DStartDisp}
  PcDc
\where
  pc~aDi = pidle \\
  dcs' = cdisp \\
  pc' = pc \oplus \{ aDi \mapsto pend \}
\end{schema}

\begin{schema}{DStartUnresp}
  PcDc
\where
  pc~aDi = pidle \\
  dcs' = cunresp \\
  pc' = pc \oplus \{ aDi \mapsto pend \}
\end{schema}

A discoverer: a write that failed with a broken pipe (the display's end, so the
display is dead), or the keepalive giving up (unresponsive, the display perhaps
alive).

\begin{schema}{RecEnd}
  DialPc
\where
  pc~aDi = pend \\
  current = fset \\
  (dcs = cunresp \lor eof = fset) \\
  slot' = slot \\
  debt' = debt \\
  current' = current \\
  sock' = sock \\
  capN' = capN \\
  hubreq' = hubreq \\
  ((ended = nocause \land ended' = dcs) \lor (ended \neq nocause \land ended' = ended)) \\
  pc' = pc \oplus \{ aDi \mapsto psu \}
\end{schema}

\begin{schema}{RecGone}
  OnlyPc
\where
  pc~aDi = pend \\
  current = fclear \\
  pc' = pc \oplus \{ aDi \mapsto pdone \}
\end{schema}

The cause is recorded, first wins.  A dial no longer current is left alone.

\begin{schema}{SPStart}
  OnlyPc
\where
  pc~aSp = pidle \\
  pc' = pc \oplus \{ aSp \mapsto phub \}
\end{schema}

The stopper: \texttt{stop} (when a listener exists) or \texttt{close}.  Both end
by asking for the Hub's end.  A \texttt{stop} on a dial with no listener
returns without ending anything; that is a subset of \texttt{close}.

\begin{schema}{HubReq}
  DialPc \\
  a? : ACTOR
\where
  a? \in \{ aSt, aSp \} \\
  pc~a? = phub \\
  current = fset \\
  ended' = ended \\
  slot' = slot \\
  debt' = debt \\
  current' = current \\
  sock' = sock \\
  capN' = capN \\
  hubreq' = fset \\
  pc' = pc \oplus \{ a? \mapsto psu \}
\end{schema}

\begin{schema}{HubGone}
  OnlyPc \\
  a? : ACTOR
\where
  a? \in \{ aSt, aSp \} \\
  pc~a? = phub \\
  current = fclear \\
  pc' = pc \oplus \{ a? \mapsto pdone \}
\end{schema}

The Hub's end sets a request.  The request itself is what refuses writes and
fails waiters at once (rule~5); it records no cause and captures nothing.  That
refusal of writes and failing of waiters is a stated rule, not modelled state:
no checked property depends on it.

\section{The no-reader settle, the wait and the release}

\begin{schema}{SUGone}
  OnlyPc \\
  a? : ACTOR
\where
  a? \in \{ aSt, aSp, aDi \} \\
  pc~a? = psu \\
  current = fclear \\
  pc' = pc \oplus \{ a? \mapsto pdone \}
\end{schema}

\begin{schema}{SUGrant}
  DialPc \\
  a? : ACTOR
\where
  a? \in \{ aSt, aSp, aDi \} \\
  pc~a? = psu \\
  current = fset \\
  slot = sfree \\
  slot' = sinline \\
  ended' = ended \\
  debt' = debt \\
  current' = current \\
  sock' = sock \\
  capN' = capN \\
  hubreq' = hubreq \\
  pc' = pc \oplus \{ a? \mapsto pset \}
\end{schema}

\begin{schema}{SURefuse}
  DialPc \\
  a? : ACTOR
\where
  a? \in \{ aSt, aSp, aDi \} \\
  pc~a? = psu \\
  current = fset \\
  slot \neq sfree \\
  debt' = fset \\
  ended' = ended \\
  slot' = slot \\
  current' = current \\
  sock' = sock \\
  capN' = capN \\
  hubreq' = hubreq \\
  pc' = pc \oplus \{ a? \mapsto pwait \}
\end{schema}

\texttt{settle\_unread}.  A free slot is taken and the caller settles.  A held
slot is left a debt, in the same step as the refusal: the holder cannot release
between the refusal and the debt.  The asker then waits.

\begin{schema}{WaitGone}
  OnlyPc \\
  a? : ACTOR
\where
  a? \in \{ aSt, aSp, aDi \} \\
  pc~a? = pwait \\
  current = fclear \\
  pc' = pc \oplus \{ a? \mapsto pdone \}
\end{schema}

\begin{schema}{Timeout}
  CapCarve \\
  a? : ACTOR
\where
  a? \in \{ aSt, aSp, aDi \} \\
  pc~a? = pwait \\
  current = fset \\
  ((ended = nocause \land ended' = chub) \lor (ended \neq nocause \land ended' = ended)) \\
  current' = fclear \\
  capN' = capN + 1 \\
  slot' = slot \\
  debt' = debt \\
  sock' = sock \\
  hubreq' = hubreq \\
  s2bad' = s2bad \\
  ((capN \geq 1 \land s3bad' = fset) \lor (capN = 0 \land s3bad' = s3bad)) \\
  s4bad' = s4bad \\
  s5bad' = s5bad \\
  ((ended' \in \{ cdisp, cunresp \} \land sock > 0 \land s2blk' = fset) \lor
    (\lnot (ended' \in \{ cdisp, cunresp \} \land sock > 0) \land s2blk' = s2blk)) \\
  s2dec' = s2dec \\
  ((ended = nocause \land eof = fset \land s5blk' = fset) \lor
    (\lnot (ended = nocause \land eof = fset) \land s5blk' = s5blk)) \\
  pc' = pc \oplus \{ a? \mapsto pdone \}
\end{schema}

Rule~6: the bound has passed.  The asker records its cause with the first-wins
record and captures, though it does not hold the slot and does not drain.  It is
the one capture that is not by a holder.  A \texttt{hub} capture charges nothing.
The two flags say what this costs: a display end captured with an
acknowledgement unread (\texttt{s2blk}), and a \texttt{hub} recorded for a
display that has died unobserved (\texttt{s5blk}).

\begin{schema}{SDrain}
  Drain \\
  a? : ACTOR
\where
  pc~a? = pset \\
  slot \neq sfree \\
  current = fset \\
  (ended \neq nocause \lor hubreq = fset \lor eof = fset) \\
  sock' = 0 \\
  ended' = ended \\
  slot' = slot \\
  debt' = debt \\
  current' = current \\
  capN' = capN \\
  hubreq' = hubreq \\
  ((eof = fset \land obs' = obs \oplus \{ a? \mapsto cdisp \}) \lor
    (eof = fclear \land obs' = obs \oplus \{ a? \mapsto nocause \})) \\
  dsaw' = dsaw \oplus \{ a? \mapsto ended \} \\
  pc' = pc \oplus \{ a? \mapsto pdec \}
\end{schema}

The drain: the socket is read to the point of blocking.  It notes whether end of
file came, and which cause stood when it began.

\begin{schema}{SDecide}
  Cap \\
  a? : ACTOR
\where
  pc~a? = pdec \\
  slot \neq sfree \\
  current = fset \\
  ((ended = nocause \land obs~a? \neq nocause \land ended' = obs~a?) \lor
    (ended = nocause \land obs~a? = nocause \land hubreq = fset \land ended' = chub) \lor
    (ended \neq nocause \land ended' = ended \land
      \lnot (ended \in \{ cdisp, cunresp \} \land dsaw~a? = nocause))) \\
  sock' = sock \\
  current' = fclear \\
  capN' = capN + 1 \\
  slot' = slot \\
  debt' = debt \\
  hubreq' = hubreq \\
  ((ended' = cdisp \land sock > 0 \land s2bad' = fset) \lor
    (\lnot (ended' = cdisp \land sock > 0) \land s2bad' = s2bad)) \\
  ((capN \geq 1 \land s3bad' = fset) \lor (capN = 0 \land s3bad' = s3bad)) \\
  s4bad' = s4bad \\
  s5bad' = s5bad \\
  pc' = pc \oplus \{ a? \mapsto nxt~a? \}
\end{schema}

The decision, by a holder: the first-wins record of the cause, then the
capture.  The cause is the end of file the drain saw, else a pending Hub
request; a cause already recorded stands, unless it is a display or
unresponsive end recorded since the drain began, in which case this step is not
enabled and \texttt{SRedrain} is.

\begin{schema}{SRedrain}
  OnlyPc \\
  a? : ACTOR
\where
  pc~a? = pdec \\
  slot \neq sfree \\
  current = fset \\
  ended \in \{ cdisp, cunresp \} \\
  dsaw~a? = nocause \\
  pc' = pc \oplus \{ a? \mapsto pset \}
\end{schema}

The record was lost to, or made after, the drain: drain again.  After this the
record stands and nothing can change the cause, so this loop runs at most once.

\begin{schema}{SettleGone}
  OnlyPc \\
  a? : ACTOR
\where
  pc~a? \in \{ pset, pdec \} \\
  current = fclear \\
  pc' = pc \oplus \{ a? \mapsto nxt~a? \}
\end{schema}

The dial was captured by someone else meanwhile (the bound of rule~6).  The
holder does nothing.

\begin{schema}{RelOwes}
  DialPc \\
  a? : ACTOR
\where
  a? \in \{ aSt, aSp, aDi, aIn \} \\
  pc~a? = prel \\
  slot = sinline \\
  debt = fset \\
  debt' = fclear \\
  ended' = ended \\
  slot' = slot \\
  current' = current \\
  sock' = sock \\
  capN' = capN \\
  hubreq' = hubreq \\
  pc' = pc \oplus \{ a? \mapsto pset \}
\end{schema}

\begin{schema}{RelRead}
  OnlyPc \\
  a? : ACTOR
\where
  a? \in \{ aSt, aSp, aDi, aIn \} \\
  pc~a? = prel \\
  slot = sinline \\
  debt = fclear \\
  pc' = pc \oplus \{ a? \mapsto prelf \}
\end{schema}

\begin{schema}{RelFree}
  DialPc \\
  a? : ACTOR
\where
  a? \in \{ aSt, aSp, aDi, aIn \} \\
  pc~a? = prelf \\
  slot = sinline \\
  slot' = sfree \\
  ended' = ended \\
  debt' = debt \\
  current' = current \\
  sock' = sock \\
  capN' = capN \\
  hubreq' = hubreq \\
  pc' = pc \oplus \{ a? \mapsto pdone \}
\end{schema}

THE DEFECT (obligation 2).  The release reads the debt (clear) in one step and
frees the slot in a later one: a refusal that lands between sets the debt on a
slot the releaser is about to free.  \texttt{RelFree} is split into
\texttt{RelRead} and this step, and \texttt{prelf} is added to \texttt{PC}.

The release.  A debt, left by a refusal that came before it, is honoured by
keeping the slot and settling again.  A debt left after the release finds the
slot free.

\section{The inline reader}

\begin{schema}{IAcq}
  DialPc
\where
  pc~aIn = pidle \\
  slot = sfree \\
  slot' = sinline \\
  ended' = ended \\
  debt' = debt \\
  current' = current \\
  sock' = sock \\
  capN' = capN \\
  hubreq' = hubreq \\
  pc' = pc \oplus \{ aIn \mapsto iread \}
\end{schema}

\begin{schema}{IReadIdle}
  OnlyPc
\where
  pc~aIn = iread \\
  blocked = fclear \\
  pc' = pc \oplus \{ aIn \mapsto ilast \}
\end{schema}

\begin{schema}{IReadNotice}
  OnlyPc
\where
  pc~aIn = iread \\
  blocked = fclear \\
  (ended \neq nocause \lor hubreq = fset \lor eof = fset) \\
  pc' = pc \oplus \{ aIn \mapsto pset \}
\end{schema}

\begin{schema}{ILastSet}
  OnlyPc
\where
  pc~aIn = ilast \\
  current = fset \\
  (ended \neq nocause \lor hubreq = fset) \\
  pc' = pc \oplus \{ aIn \mapsto pset \}
\end{schema}

\begin{schema}{ILastSkip}
  OnlyPc
\where
  pc~aIn = ilast \\
  (current = fclear \lor (ended = nocause \land hubreq = fclear)) \\
  pc' = pc \oplus \{ aIn \mapsto prel \}
\end{schema}

\begin{schema}{IBad}
  CapCarve
\where
  pc~aIn = iread \\
  blocked = fclear \\
  current = fset \\
  ((ended = nocause \land ended' = cdisp) \lor (ended \neq nocause \land ended' = ended)) \\
  current' = fclear \\
  capN' = capN + 1 \\
  slot' = slot \\
  debt' = debt \\
  sock' = sock \\
  hubreq' = hubreq \\
  s2bad' = s2bad \\
  ((capN \geq 1 \land s3bad' = fset) \lor (capN = 0 \land s3bad' = s3bad)) \\
  s4bad' = s4bad \\
  s5bad' = s5bad \\
  s2blk' = s2blk \\
  ((ended' = cdisp \land sock > 0 \land s2dec' = fset) \lor
    (\lnot (ended' = cdisp \land sock > 0) \land s2dec' = s2dec)) \\
  s5blk' = s5blk \\
  pc' = pc \oplus \{ aIn \mapsto prel \}
\end{schema}

A caller that holds the slot reads, notices an end (or a request) and owes a
settle, or finishes its read; its last check, \texttt{\_settle\_if\_ended}, then
either finds an end and owes a settle or finds none and releases.  Between the
last check and the release, anyone may end the dial: that is the window the
debt closes.  A frame that breaks the stream (\texttt{IBad}) records the
display's end and captures at once, undrained: the first carve-out.

\section{The listener thread}

\begin{schema}{RNotice}
  OnlyPc
\where
  pc~aRd = rtop \\
  blocked = fclear \\
  (hubreq = fset \lor ended \neq nocause \lor eof = fset) \\
  pc' = pc \oplus \{ aRd \mapsto pset \}
\end{schema}

\begin{schema}{RBad}
  CapCarve
\where
  pc~aRd = rtop \\
  blocked = fclear \\
  current = fset \\
  ((ended = nocause \land ended' = cdisp) \lor (ended \neq nocause \land ended' = ended)) \\
  current' = fclear \\
  capN' = capN + 1 \\
  slot' = slot \\
  debt' = debt \\
  sock' = sock \\
  hubreq' = hubreq \\
  s2bad' = s2bad \\
  ((capN \geq 1 \land s3bad' = fset) \lor (capN = 0 \land s3bad' = s3bad)) \\
  s4bad' = s4bad \\
  s5bad' = s5bad \\
  s2blk' = s2blk \\
  ((ended' = cdisp \land sock > 0 \land s2dec' = fset) \lor
    (\lnot (ended' = cdisp \land sock > 0) \land s2dec' = s2dec)) \\
  s5blk' = s5blk \\
  pc' = pc \oplus \{ aRd \mapsto rdone \}
\end{schema}

The listener's loop tests, in turn, the Hub's request (its stop flag), a
recorded end, and the display's end of file, and on any of them settles through
the same two steps that every other holder uses.  It records no cause of its
own and reads none twice.  A malformed frame is the first carve-out
(\texttt{RBad}).

\section{Refinement of crash\_attribution\_inflight\_ack\_drain.tex}

The abstract model that this specification refines is
\texttt{crash\_attribution\_inflight\_ack\_drain.tex}, the variant of
\texttt{crash\_attribution.tex} in which every \texttt{Discover*} step takes each
acknowledgement in the socket and in the reader's buffer into \texttt{acked} in
the same step (\texttt{Recv} and \texttt{Deliver} folded into the capture), and
in which a drain that applies an ack clears the carried set.  Its \texttt{ended}
takes \texttt{edisp}, \texttt{eunr} or \texttt{ehub}.  The gluing, concrete to
abstract:

\begin{itemize}
  \item \textbf{Cause.}  \texttt{cdisp}, \texttt{cunresp}, \texttt{chub} map to
    \texttt{edisp}, \texttt{eunr}, \texttt{ehub}.  A capture here is the abstract
    discovery of that cause.  The concrete \texttt{hubreq} is an abstract
    stutter: no abstract state changes when the Hub merely asks.  The abstract
    \texttt{HubEnd} is placed at the settle's observation (the drain), not at the
    ask.  That placement is the linearisation that makes the S5 residual a
    display that died \emph{after} the abstract \texttt{HubEnd}, that is before
    the abstract \texttt{CrashFrame}.
  \item \textbf{The drain.}  The concrete guard ``no ack unread when
    \texttt{cdisp} is captured'' (\texttt{sock = 0} at \texttt{SDecide}) refines
    the abstract in-step drain of the variant.  \texttt{sock} conflates the socket
    and the reader's undelivered buffer.  That is sound only because rule~1 makes
    the settler the owner of the frame buffer: a non-owner cannot empty what the
    reader has taken, but the owner's drain takes it to the ledger before it
    returns, so for the owner the two are one place.  This is the strongest
    argument for rule~1.  The abstract \texttt{Recv} then \texttt{Deliver} is one
    concrete step because the settler runs both without releasing the slot.
  \item \textbf{The gluing is an inclusion, not an equality of state.}  The
    unresponsive re-drain reads acknowledgements that the abstract model says
    were never sent: after \texttt{Unresponsive} a \texttt{Swap} adds nothing to
    the abstract \texttt{flight}.  The concrete \texttt{acked} is therefore a
    superset of the abstract one, and the relation that holds is
    \emph{concrete charges $\subseteq$ abstract charges}.  The inclusion gives A2.
    A1 holds for a different reason: the only acknowledgements the concrete
    settle reads beyond the abstract ones are, by T11, of scenes that survived a
    frame, so the poison scene is never among them.  The abstract model charges
    a scene that the concrete one does not only when a display recorded
    unresponsive later swaps and acknowledges that scene, and such a scene was
    not poison.
  \item \textbf{The obligations.}  A3 (a Hub-ended dial charges nothing) holds by
    the cause mapping, and by S3 (one capture).  A4 (a death is charged once)
    holds by S3.  A1 (an observed crash is charged) needs S1 and S5.  A2 (no
    scene outside the ledger is charged) needs S2.  O2's drain clause is S2.  All
    five hold on the holder's settle, with the re-drain of rule~4; without it A2
    and O2 fail (\texttt{split\_settle}).
\end{itemize}

\paragraph{Four exceptions to A1 and A2.}
The carve-outs are not bounds that the abstract model already contains.
\texttt{crash\_attribution.tex} has no charging stream-error end: its reply skew
is \texttt{HubEnd}, and every display or unresponsive discovery is guarded on an
empty flight and an empty reader buffer, so it never charges a scene whose ack
is unread.  Four concrete paths have no abstract counterpart.  On each, A1 or A2
fails by design; they are deliberate weakenings, not refinements.
\begin{enumerate}
  \item \texttt{s2dec}: a display charge with an ack unread, from a frame that
    broke the stream.  A2 fails.  Bound: the acks behind the bad frame, once.
  \item \texttt{s2blk}, display cause: a display charge with an ack unread,
    taken by the bound of rule~6.  A2 fails.  Bound: every ack queued while the
    holder stalled, including the holder's undelivered buffer, after the stop
    join.  The over-charge is accepted (DES-102).
  \item \texttt{s2blk}, unresponsive cause: the bound's capture of an
    unresponsive end, undrained.  A2 fails.  The bound is that of item~2.  The
    abstract model's B3 makes the retire instant the capture, but it still
    guards the charge on an empty flight and buffer.
  \item \texttt{s5blk}: \texttt{hub} recorded by the bound for a display that has
    already died.  An abstract \texttt{HubEnd} is disabled once the display is
    down.  The concrete result is that a poison death is carried and not charged,
    and a later ack read clears it.  A1 is lost by the B4 route, once, and only
    after the bound.
\end{enumerate}

This section is the refinement obligation list.  The matching note belongs in
\texttt{crash\_attribution.tex} next to its settle, and in the Bounds of DES-102;
the first is outside this mission's write set, so it is not edited here.

\section{Verification}

Commands (each run alone; MAX\_INITIALISATIONS must be at least 4: the slot
born with the listener or free, an unread acknowledgement or none).  Search mode
is stated with every count, because the breadth-first count and the default
count differ.

\begin{verbatim}
  fuzz docs/dial_settlement.tex                       (clean)
  P="-p MAX_OPERATIONS 100000 -p MAX_INITIALISATIONS 100"
  probcli docs/dial_settlement.tex -model_check -coverage $P          (deadlock, default mode)
  probcli docs/dial_settlement.tex -model_check -nodead -bf $P \
    -goal "s2bad = fset or s3bad = fset or s4bad = fset or s5bad = fset"
  S1a (timer-free; pwait counts as waiting, a stalled holder does not count):
  probcli docs/dial_settlement.tex -model_check -nodead -bf $P \
    -goal "current = fset & (ended /= nocause or hubreq = fset) &
           !a.(a : ACTOR => pc(a) : {pidle,pdone,rdone,rnone,pwait})"
  S1b (a stalled holder counts as stuck; a waiting asker is not):
  probcli docs/dial_settlement.tex -model_check -nodead -bf $P \
    -goal "current = fset & (ended /= nocause or hubreq = fset) &
           !a.(a : ACTOR => (pc(a) : {pidle,pdone,rdone,rnone} or
                (blocked = fset & pc(a) : {rtop,iread})))"
  The carve-outs are reachable, one flag each:
  probcli ... -goal "s2blk = fset"   / "s2dec = fset"   / "s5blk = fset"
\end{verbatim}

\paragraph{Results for this control.}
S1a goal, \texttt{-nodead -bf}: found, 5932 states, 8 steps.  The release reads
no debt, a refusal sets the debt, and the release then frees the slot.

\end{document}
